% GENERATED FILE. Edit canonical lessons, then run npm run generate:books.
% canonical-source-hash: fnv1a64:21f786f60b9c65ca
% canonical-lessons: AR-C507-marida, AR-C507-nazlat-bard, AR-C507-ghathayan, AR-C507-taqayyaa, AR-C507-ishal

\chapter{Falling ill}
\label{ch:ar-falling-ill}
\hlchaptermodality{507}

\begin{chapteropening}
\textbf{By the end of this chapter:} \emph{I can say that someone fell ill or vomited, and talk about a chill, nausea and diarrhoea.}

Complete the last lesson of chapter 507: I can say that someone fell ill or vomited, and talk about a chill, nausea and diarrhoea.
\end{chapteropening}

\section[\texorpdfstring{\ar{مرض} (mariḍa)}{مرض (mariḍa)}]{\ar{مرض} (mariḍa) — to fall ill}

\label{lesson:AR-C507-marida}



\begin{quote}
Before the new one: say the Arabic for a chin, then the Arabic for a nerve.
\end{quote}

\subsection*{You'll want to know: \ar{مرض}}
\textbf{\ar{مرض}} — \emph{mariḍa} — \textquotedblleft{}to fall ill\textquotedblright{}.

\begin{quote}
\ar{مرض} \ar{أخي} \ar{في} \ar{الشتاء} \ar{الماضي}.
\end{quote}

\begin{quote}
\emph{mariḍa akhī fī ash-shitāʾ al-māḍī.}
\end{quote}

\begin{quote}
My brother fell ill last winter.
\end{quote}

Present: \ar{يمرض}, \emph{yamraḍ}.

The noun \ar{مرض}, read maraḍ, is an illness.

\begin{cousinweb}
\textbf{\ar{مرض}} (\emph{mariḍa}) is built on the native Arabic root \emph{m-r-ḍ}, \textquotedblleft{}to be sick\textquotedblright{} (Lane). Marīḍ, sick, which you know, is from the same root.
\end{cousinweb}

\begin{grammarlens}[title={What you've built}]
One more word for describing what you lived and felt.
\end{grammarlens}

\subsection*{Your turn}
\begin{itemize}
\raggedright
  \item \emph{Say it:} \emph{mariḍa}
  \item \emph{Say it:} \emph{mariḍa}, once more
  \item \emph{Say it:} \emph{mariḍa akhī fī ash-shitāʾ al-māḍī.} — My brother fell ill last winter.
  \item \emph{Recall:} say the Arabic for a chin, then the Arabic for a nerve, then say \emph{mariḍa} again
\end{itemize}

\begin{tcolorbox}[breakable,colback=teal!4,colframe=teal!35!black,title={Before you move on}]
What does \ar{مرض} mean? (\textquotedblleft{}To fall ill\textquotedblright{}.) Say it once more.
\end{tcolorbox}
\section[\texorpdfstring{\ar{نزلة} \ar{برد} (nazlat bard)}{نزلة برد (nazlat bard)}]{\ar{نزلة} \ar{برد} (nazlat bard) — a chill}

\label{lesson:AR-C507-nazlat-bard}



\begin{quote}
Before the new one: say the Arabic for a nerve, then the Arabic for to fall ill.
\end{quote}

\subsection*{You'll want to know: \ar{نزلة} \ar{برد}}
\textbf{\ar{نزلة} \ar{برد}} — \emph{nazlat bard} — \textquotedblleft{}a chill\textquotedblright{}.

\begin{quote}
\ar{عند} \ar{أخي} \ar{نزلة} \ar{برد} \ar{خفيفة}.
\end{quote}

\begin{quote}
\emph{ʿinda akhī nazlat bard khafīfa.}
\end{quote}

\begin{quote}
My brother has a slight chill.
\end{quote}

\begin{cousinweb}
\textbf{\ar{نزلة} \ar{برد}} (\emph{nazlat bard}) is built on two native Arabic roots: \emph{nazla}, \textquotedblleft{}a coming down\textquotedblright{}, from \emph{n-z-l}, the root of \emph{nazala}, to go down (Lane), and \emph{bard}, \textquotedblleft{}cold\textquotedblright{}: a cold that comes down on you.
\end{cousinweb}

\begin{grammarlens}[title={What you've built}]
One more word for describing what you lived and felt.
\end{grammarlens}

\subsection*{Your turn}
\begin{itemize}
\raggedright
  \item \emph{Say it:} \emph{nazlat bard}
  \item \emph{Say it:} \emph{nazlat bard}, once more
  \item \emph{Say it:} \emph{ʿinda akhī nazlat bard khafīfa.} — My brother has a slight chill.
  \item \emph{Recall:} say the Arabic for a nerve, then the Arabic for to fall ill, then say \emph{nazlat bard} again
\end{itemize}

\begin{tcolorbox}[breakable,colback=teal!4,colframe=teal!35!black,title={Before you move on}]
What does \ar{نزلة} \ar{برد} mean? (\textquotedblleft{}A chill\textquotedblright{}.) Say it once more.
\end{tcolorbox}
\section[\texorpdfstring{\ar{غثيان} (ghathayān)}{غثيان (ghathayān)}]{\ar{غثيان} (ghathayān) — nausea}

\label{lesson:AR-C507-ghathayan}



\begin{quote}
Before the new one: say the Arabic for to fall ill, then the Arabic for a chill.
\end{quote}

\subsection*{You'll want to know: \ar{غثيان}}
\textbf{\ar{غثيان}} — \emph{ghathayān} — \textquotedblleft{}nausea\textquotedblright{}.

\begin{quote}
\ar{شعرت} \ar{بالغثيان} \ar{في} \ar{السفينة}.
\end{quote}

\begin{quote}
\emph{shaʿartu bi-l-ghathayān fī as-safīna.}
\end{quote}

\begin{quote}
I felt sick on the ship.
\end{quote}

\begin{cousinweb}
\textbf{\ar{غثيان}} (\emph{ghathayān}) is built on the native Arabic root \emph{gh-th-y}, \textquotedblleft{}to be troubled, to heave (of the stomach)\textquotedblright{} (Lane). The pattern faʿalān often names a restless movement.
\end{cousinweb}

\begin{grammarlens}[title={What you've built}]
One more word for describing what you lived and felt.
\end{grammarlens}

\subsection*{Your turn}
\begin{itemize}
\raggedright
  \item \emph{Say it:} \emph{ghathayān}
  \item \emph{Say it:} \emph{ghathayān}, once more
  \item \emph{Say it:} \emph{shaʿartu bi-l-ghathayān fī as-safīna.} — I felt sick on the ship.
  \item \emph{Recall:} say the Arabic for to fall ill, then the Arabic for a chill, then say \emph{ghathayān} again
\end{itemize}

\begin{tcolorbox}[breakable,colback=teal!4,colframe=teal!35!black,title={Before you move on}]
What does \ar{غثيان} mean? (\textquotedblleft{}Nausea\textquotedblright{}.) Say it once more.
\end{tcolorbox}
\section[\texorpdfstring{\ar{تقيأ} (taqayyaʾa)}{تقيأ (taqayyaʾa)}]{\ar{تقيأ} (taqayyaʾa) — to vomit, to be sick}

\label{lesson:AR-C507-taqayyaa}



\begin{quote}
Before the new one: say the Arabic for a chill, then the Arabic for nausea.
\end{quote}

\subsection*{You'll want to know: \ar{تقيأ}}
\textbf{\ar{تقيأ}} — \emph{taqayyaʾa} — \textquotedblleft{}to vomit, to be sick\textquotedblright{}.

\begin{quote}
\ar{تقيأ} \ar{الطفل} \ar{بعد} \ar{الرحلة} \ar{الطويلة}.
\end{quote}

\begin{quote}
\emph{taqayyaʾa aṭ-ṭifl baʿda ar-riḥla aṭ-ṭawīla.}
\end{quote}

\begin{quote}
The child was sick after the long journey.
\end{quote}

Present: \ar{يتقيأ}, \emph{yataqayyaʾ}.

\begin{cousinweb}
\textbf{\ar{تقيأ}} (\emph{taqayyaʾa}) is built on the native Arabic root \emph{q-y-ʾ}, \textquotedblleft{}to vomit\textquotedblright{} (Lane). Form V taqayyaʾa is the usual verb today.
\end{cousinweb}

\begin{grammarlens}[title={What you've built}]
One more word for describing what you lived and felt.
\end{grammarlens}

\subsection*{Your turn}
\begin{itemize}
\raggedright
  \item \emph{Say it:} \emph{taqayyaʾa}
  \item \emph{Say it:} \emph{taqayyaʾa}, once more
  \item \emph{Say it:} \emph{taqayyaʾa aṭ-ṭifl baʿda ar-riḥla aṭ-ṭawīla.} — The child was sick after the long journey.
  \item \emph{Recall:} say the Arabic for a chill, then the Arabic for nausea, then say \emph{taqayyaʾa} again
\end{itemize}

\begin{tcolorbox}[breakable,colback=teal!4,colframe=teal!35!black,title={Before you move on}]
What does \ar{تقيأ} mean? (\textquotedblleft{}To vomit, to be sick\textquotedblright{}.) Say it once more.
\end{tcolorbox}
\section[\texorpdfstring{\ar{إسهال} (ishāl)}{إسهال (ishāl)}]{\ar{إسهال} (ishāl) — diarrhoea}

\label{lesson:AR-C507-ishal}



\begin{quote}
Before the new one: say the Arabic for nausea, then the Arabic for to vomit.
\end{quote}

\subsection*{You'll want to know: \ar{إسهال}}
\textbf{\ar{إسهال}} — \emph{ishāl} — \textquotedblleft{}diarrhoea\textquotedblright{}.

\begin{quote}
\ar{أعطى} \ar{الطبيب} \ar{الطفل} \ar{دواء} \ar{للإسهال}.
\end{quote}

\begin{quote}
\emph{aʿṭā aṭ-ṭabīb aṭ-ṭifl dawāʾ li-l-ishāl.}
\end{quote}

\begin{quote}
The doctor gave the child medicine for diarrhoea.
\end{quote}

\begin{cousinweb}
\textbf{\ar{إسهال}} (\emph{ishāl}) is built on the native Arabic root \emph{s-h-l}, \textquotedblleft{}to be smooth, easy\textquotedblright{} (Lane). Lane gives Form IV ashala, of a medicine, as to loosen the bowels; sahl, easy, is the same root.
\end{cousinweb}

\begin{grammarlens}[title={What you've built}]
One more word for describing what you lived and felt.
\end{grammarlens}

\subsection*{Your turn}
\begin{itemize}
\raggedright
  \item \emph{Say it:} \emph{ishāl}
  \item \emph{Say it:} \emph{ishāl}, once more
  \item \emph{Say it:} \emph{aʿṭā aṭ-ṭabīb aṭ-ṭifl dawāʾ li-l-ishāl.} — The doctor gave the child medicine for diarrhoea.
  \item \emph{Recall:} say the Arabic for nausea, then the Arabic for to vomit, then say \emph{ishāl} again
\end{itemize}

\begin{tcolorbox}[breakable,colback=teal!4,colframe=teal!35!black,title={Before you move on}]
What does \ar{إسهال} mean? (\textquotedblleft{}Diarrhoea\textquotedblright{}.) Say it once more.
\end{tcolorbox}
